\chapter{Flow-based models}

\textbf{NB:} This is based on \cite{papamakarios2021normalizing}, \cite{lipman2023flowmatching}, \cite{grathwohl2019ffjord},  \cite{tong2023improving},  \cite{chen2018neuralode}, and \cite{liu2022rectified}.


\paragraph{Motivation.}

We aim to define a distribution over $\x\in\R^D$. Rather than parametrising it directly, we will express $\x$ as a transformation of a random vector $\u$. That is,
\[\x = T(\u), \qquad \u \sim p_\u(\u).\] If the map and the base (or source) measure have parameters, say, $\phi$ and $\psi$ respectively, they will be the parameters of the target measure denoted $p(\x)$.

Formally, the induced model will be the pushforward measure of $p_\u$ through $T$, that is,
\begin{equation}
    \label{eq:push_forward_density}
    \int_A p(\x)\, d\x
=
\int_{T^{-1}(A)} p_\u(\u)\, d\u ,
\qquad
\forall A \subset \mathbb{R}^D \text{ measurable}.    
\end{equation}


This is somehow an extension of the setting we adopted in the variational autoencoder formulation, where the base distribution was assumed to be Gaussian, the transformation $T(\cdot)$ a neural network, and the learning paradigm was variational inference. The extensions to be considered in this chapter will take two directions. First, by imposing continuity and invertibility conditions of the map, the resulting model can be trained via maximum likelihood, provided the inverse of the map can be computed. Second, we will consider an approach where the map is broken down in stages, initially considering a compositional construction of the map and then an a continuous version; we will see that this perspective calls for an entirely new learning strategy. 


\section{Normalising flows}

We will now focus on a case characterised by a clear (and arguably restrictive) condition on $T$: we will require it to be a diffeomorphism. A definition is given as follows as a reminder.

\begin{definition}[Diffeomorphism]
Let $T:\mathbb{R}^D \to \mathbb{R}^D$ be a map. We say that $T$ is a \emph{diffeomorphism} if it is bijective, continuously differentiable, and its inverse $T^{-1}$ is also continuously differentiable. Equivalently, $T$ is a $\cC^1$ bijection with a $\cC^1$ inverse. In particular, the Jacobian matrix $\nabla T(\u)$ exists everywhere and is non-singular, i.e.\ $\det \nabla T(\u) \neq 0$ for all $\u \in \mathbb{R}^D$.
\end{definition}

The comment about $T$ being a bijective map with a well defined Jacobian is crucial in this setting, as that enables to  us to express the induced density over $\x = T(\u)$ via the change-of-variables theorem in terms of the base measure $p_\u$ as follows 
\begin{equation}
\label{eq:CoV_invT}    
p(\x)
=
p_\u\!\left(T^{-1}(\x)\right)
\left|
\det \nabla T^{-1}(\x)
\right|,
\end{equation}
or equivalently, as a function of the source variable $\u$:
\begin{equation}
\label{eq:CoV_T}    
p(\x)
=
p_\u(\u)
\left|
\det \nabla T(\u)
\right|^{-1}.
\end{equation}


Some relevant remarks mentioned as follows. 

\begin{remark}[Mass-preserving warping]
The map $T$ can be interpreted as a warping of the ambient space $\cX$, which reorganises the probability mass of $p_\u$ into the target density $p$. Since probability mass must be preserved, the transformation only \emph{transports} mass; it neither creates nor destroys it. Therefore, local changes in volume induced by the map must be reflected in the resulting density, and this is precisely the role of the Jacobian factor $\det \nabla T^{-1}(\x)$ in eq.~\eqref{eq:CoV_invT} above. If the transformation locally compesses the space, the density must increase accordingly to preserve mass; conversely, if the space is locally expanded, the density must decrease.
\end{remark}


\begin{remark}[Training: Closed form likelihood]
The identity in eq.~\eqref{eq:CoV_invT} allows for the likelihood, given observations $\x$, to be calculated in closed form, provided that both the inverse map $T^{-1}$ and the determinant of its Jacobian can be evaluated.
\end{remark}

\begin{remark}[Sampling]
Under the proposed methodology, generating samples from $\x\sim p(\x)$ is direct by first sampling $\u\sim p_\u$, and then evaluating $\x = T(\u)$. This can be done efficiently if the source measure is a simple density (e.g., Gaussian or uniform), and if the map $T$ can be evaluated.
\end{remark}

\begin{remark}[Normalising flow]
Since in either case both the direct map $T:\u\to\x$ when the aim is sampling, or the reverse one $T^{-1}:\x\to\u$ when the aim is learning. Since the case is usually the latter, the inverse map is usually the parametrised via $\phi=T^{-1}$, and the base measure is considered to be Gaussian. Therefore, the role of $\phi$ is to convert an arbitrary data distribution $p(\x)$ into a Normal variable, hence the name for $\phi$ as normalising flow.
\end{remark}




\paragraph{Construction of expressive flows.}

At this stage, it is natural to ask how expressive the normalising flow construction is. In other words, given an arbitrary pair of densities $p_\u(\u), p(\x)$, does a map $T$ pushing the former to the latter always exist? For the scalar case, i.e., $\x\in\R$, verifying this is straightforward: when the densities are smooth, we can consider the mononote rearrangement given by 
\begin{equation}
    T(\u) = F^{-1}_\x(F_\u(\u)),
\end{equation}
where $F_\x$ and $F_\u$ are the cumulative distribution functions of $\x$ and $\u$ respectively. Intuitively, this flow is pushing the mass from $p_\u$ to a uniform distribution using $F_\u(\cdot)$, and then pushing it from there to $p$ via $F^{-1}_\x(\cdot)$. 

Regarding the non-scalar case, provided that $p_\u$ is absolutely continuous wrt the Lebesgue measure a map fulfilling eq.~\eqref{eq:push_forward_density} is guaranteed to exist based on Brenier Theorem (see the OT chapter). this result states that the OT between two a.c. measures under the Euclidean ground cost always exist, is unique, and (critical in the normalising flow setting) is the gradient of a convex potential (which ensures invertibility). 

In addition to Brenier's theoretical guarantees, further insight into the expressivity of normalising flows can be gained by trying to adapt the monotone arrangement concept to the vector-valued case. Let us first observe that the monotone arrangement cannot be applied to vectors because, although the multivariate CDF is well defined and given by 
\begin{equation}
    F_\x(\x) = \int_{(-\infty,x_1)\times\cdots\times (-\infty,x_D)}p(\x)\d\x,
\end{equation}
it is not invertible and, thus, a monotone arrangement cannot be constructed (there is, in particular, no concept of motononicity in $\R^D$).

An alternative is to consider the triangular flow given by the \emph{Knothe--Rosenblatt rearrangement}---see \cite[Ch.~2]{villani2009optimal}. In this construction the
transport map is triangular, meaning that each coordinate depends
only on the current and previous ones,
\[
T(\u)
=
\begin{pmatrix}
T_1(u_1)\\
T_2(u_1,u_2)\\
\vdots\\
T_D(u_1,\ldots,u_D)
\end{pmatrix}.
\]

The map is constructed recursively using conditional CDFs. Writing
the base density as
\[
p_\u(\u)
=
p_{u,1}(u_1)
\prod_{k=2}^D
p_{u,k}(u_k \mid u_{1:k-1}),
\]
define the conditional CDFs
\[
F_{u,k}(t \mid u_{1:k-1})
=
\int_{-\infty}^{t}
p_{u,k}(s \mid u_{1:k-1})\,ds,
\]
and similarly for the target density $p(\x)$,
\[
F_k(t \mid x_{1:k-1})
=
\int_{-\infty}^{t}
p_k(s \mid x_{1:k-1})\,ds .
\]

The transport map is then defined recursively as
\[
x_k
=
F_k^{-1}\!\left(
F_{u,k}(u_k \mid u_{1:k-1})
\mid x_{1:k-1}
\right),
\qquad k=1,\ldots,D .
\]

Intuitively, each step first maps $u_k$ to a uniform variable via the
conditional CDF of the base distribution and then maps it to the
correct conditional distribution of the target via the inverse
conditional CDF. This construction yields a diffeomorphic triangular
map $T$ satisfying $T_\# p_\u = p$, showing that flows are expressive
enough to represent arbitrary smooth densities.

This triangular structure is particularly attractive from the
perspective of the change-of-variables formula. Indeed, the Jacobian
matrix $\nabla T(\u)$ is triangular, and therefore its determinant is
simply the product of its diagonal entries,
\[
\det \nabla T(\u)
=
\prod_{k=1}^D
\frac{\partial T_k}{\partial u_k}(u_{1:k}).
\]
As a result, the Jacobian determinant required in
eq.~\eqref{eq:CoV_T} can be computed efficiently, which facilitates training. Additionally, since the inverse of a triangular map is also  triangular, the same computational advantage holds when evaluating the likelihood
via the inverse transformation in eq.~\eqref{eq:CoV_invT}.


Another natural way to construct expressive flows is to build the map as a
composition of simpler diffeomorphisms. Concretely, write
\[
T = T_K \circ \cdots \circ T_1,
\]
where each $T_k : \R^D \to \R^D$ is an invertible transformation with a
tractable Jacobian determinant. If we define the intermediate variables
$\z_0=\u$, $\z_k = T_k(\z_{k-1})$ for $k=1,\ldots,K$, and $\z_K=\x$, then the
change-of-variables formula can be applied recursively, yielding
\[
\log p(\x)
=
\log p_\u(\u)
-
\sum_{k=1}^K
\log \left|\det \nabla T_k(\z_{k-1})\right|.
\]
This compositional construction allows complex transformations to be
obtained by stacking simple building blocks, while keeping likelihood
evaluation tractable.


\missing{Include a couple of examples, perhaps using GPs}

\section{Continuous-time flows}

We now consider a continuous version of the flow, where rather than a collection of sequential transformations as in the compositional construction, we will implement flows that model the (infinitesimal) evolution of the variables. To this end, let us consider a stochastic process $z_t$ indexed by the temporal variable $t\in[0,1]$, such that 
\[\z_0 = \u,
\qquad
z_1 = \x.
\]

Unlike discrete flows, where the transformation is applied from one variable to the next, we now model the temporal evolution of $\z_t$, that is, 
\begin{equation}
    \label{eq:ODE}
    \frac{\d\z_t}{\dt} = g_\phi(t,\z_t).   
\end{equation}
 Here, we will require that the function $g_\phi: [0,1] \times \R^D \to \R^D$ is uniformly Lipschitz continuous in $\z$ and continuous in $t$, so that the ODE admits a unique solution. Furthermore, we will refer to the function $g_\phi(\cdot,\cdot)$ as \emph{velocity field}.

\begin{remark}[Neural ordinary differential equations]
The dynamics above can be interpreted as a Neural ODE. To see this, recall that a deep residual network updates a hidden state through layers of the form
\[
\z_{k+1} = \z_k + h\, g_\phi(t_k, \z_k),
\]
where $h$ is a small step size and $g_\phi$ is typically implemented as a neural network. If we interpret the layer index as a discretised time variable $t_k = kh$, this recursion corresponds to a forward
Euler discretisation of the differential equation
\[
\frac{\d \z_t}{\d t} = g_\phi(t,\z_t).
\]
In the limit, as the number of layers increases and the step size
$h\to0$, the residual network converges to the continuous-time
dynamics above. Neural ODEs \cite{chen2018neuralode} therefore provide a continuous-depth analogue of deep networks, where the transformation between input and output is defined implicitly as the solution of a learned dynamical
system.
\end{remark}

\paragraph{Evolution of densities.}
The ODE in  eq.~\eqref{eq:ODE} describes how individual samples evolve. A natural question, however, is how the corresponding probability density evolves over time based on those dynamics. Let $p_t(\z)$ denote the density of $\z_t$. Since the flow
simply transports probability mass without creating or destroying it,
the density must satisfy the \emph{continuity equation} \cite[Ch.~5]{villani2009optimal}
\begin{equation}
\label{eq:continuity}
\partial_t p_t(\z)
+
\nabla_\z \cdot \big(p_t(\z)\, g_\phi(t,\z)\big)
=
0.
\end{equation}
This equation expresses conservation of probability mass: the change in
density at a point is determined by the divergence of the probability
flux $p_t(\z) g_\phi(t,\z)$.

While eq.~\eqref{eq:continuity} describes the evolution of the full density field, understanding the change of the density along individual trajectories of the flow is insightful. To this end, we will let $\z_t$ follow the dynamics above and inspect the time derivative of the log-density along this trajectory.

Expanding the divergence in eq.~\eqref{eq:continuity} using the product
rule yields
\[
\partial_t p_t(\z)
+
g_\phi(t,\z)\cdot\nabla_\z p_t(\z)
+
p_t(\z)\,\nabla_\z\cdot g_\phi(t,\z)
=
0.
\]
Now consider a trajectory $\z_t$ satisfying
\[
\frac{\d \z_t}{\d t} = g_\phi(t,\z_t).
\]
By the chain rule, the total derivative of the density along the
trajectory is
\[
\frac{\d}{\d t} p_t(\z_t)
=
\partial_t p_t(\z_t)
+
g_\phi(t,\z_t)\cdot\nabla_\z p_t(\z_t).
\]
Substituting this expression into the expanded continuity equation
gives
\[
\frac{\d}{\d t} p_t(\z_t)
+
p_t(\z_t)\,\nabla_\z\cdot g_\phi(t,\z_t)
=
0.
\]
Dividing by $p_t(\z_t)$ yields
\[
\frac{\d}{\d t} \log p_t(\z_t)
=
-\,\nabla_\z\cdot g_\phi(t,\z_t).
\]
Finally, recalling that the divergence of a vector field equals the
trace of its Jacobian,
\[
\nabla_\z\cdot g_\phi(t,\z)
=
\mathrm{Tr}\!\left(
\frac{\partial g_\phi}{\partial \z}(t,\z)
\right),
\]
we obtain
\[
\frac{\d}{\d t} \log p_t(\z_t)
=
- \mathrm{Tr}\!\left(
\frac{\partial g_\phi}{\partial \z}(t,\z_t)
\right).
\]

\begin{remark}[Instantaneous change-of-variables formula]
The above identity is known as the \emph{instantaneous change-of-variables} formula. It can be interpreted as the continuous-time analogue of the Jacobian determinant appearing in discrete normalising flows: the divergence of the velocity field $g_\phi$ measures the instantaneous expansion or contraction of volume elements along the flow.
\end{remark}
 


\paragraph{Continuous vs discrete normalising flows.}
The formulation above highlights the appeal of continuous normalising flows. First, rather than specifying a finite composition of transformations $T = T_K \circ \cdots \circ T_1$, the flow is defined through the solution of a differential equation, which can be interpreted as the limit of infinitely many infinitesimal transformations. In practice, the effective number of transformations is determined adaptively by the numerical ODE solver, rather than being fixed a priori by the architecture.

Second, continuous flows avoid the need to explicitly construct inverse maps. Both sampling and density evaluation reduce to solving an initial value problem for the state $\z_t$ and the log-density evolution, removing the structural constraints required in discrete flows to ensure the existence (and tractability) of inverses.

These advantages come with some drawbacks. Solving the ODE typically requires many evaluations of the vector field $g_\phi$, making training considerably more expensive than in discrete flows. Moreover, the computation of the divergence term $\mathrm{Tr}(\partial g_\phi/\partial \z)$ can be costly in high-dimensional settings and often requires stochastic trace estimators. Finally, the accuracy of the model depends on the numerical integration scheme, introducing additional computational overhead and potential numerical error.

\begin{remark}[Free-form Jacobian of Reversible Dynamics (FFJORD)]
A major practical difficulty in continuous normalising flows is the computation of the divergence term $\mathrm{Tr}(\partial g_\phi/\partial \z)$ appearing in the instantaneous change-of-variables formula. In high dimensions, explicitly forming the Jacobian is prohibitively expensive. The FFJORD method \cite{grathwohl2019ffjord} addresses this issue by estimating the trace using Hutchinson's stochastic trace estimator,
\[
\mathrm{Tr}(A) = \mathbb{E}_{\epsilon}[\epsilon^\top A \epsilon],
\]
which allows the divergence to be computed using Jacobian--vector products obtainable through automatic differentiation. This enables the use of unrestricted neural vector fields $g_\phi$ while keeping the computational cost manageable.
\end{remark}

We next review an alternative formulation to specify an appropriate flow $\phi$, by directly learning the velocity field through matching it to a prescribed probability flow between simple and target distributions.

\section{Flow matching}

Rather than defining the transformation through an explicit, parametric, map $T$ or through the solution of a continuous normalising flow, the flow matching perspective first specifies a so-called probability path connecting a
simple base distribution to the target distribution. Let
$\{p_t\}_{t\in[0,1]}$ 
denote a family of densities such that
\[
p_0(\z) = p_\u(\z), 
\qquad
p_1(\z) = p(\z),
\]
which interpolates between the base and target measures.

To realise such a path, consider a time-dependent flow map
$\phi_t:\R^D \to \R^D$ 
that transports samples over time. As in the previous section, the flow $\phi$ is defined as the solution of the ODE
\[
\frac{\d}{\d t}\phi_t(\u)
=
v(t,\phi_t(\u)),
\qquad
\phi_0(\u)=\u,
\]
where $v:[0,1]\times\R^D\to\R^D$ is a time-dependent velocity field. For a sample $\u\sim p_\u$, the trajectory $\z_t=\phi_t(\u)$ describes the evolution of the particle under the velocity field $v$, and the distribution of $\z_t$ evolves along the prescribed probability path
$\{p_t\}_{t\in[0,1]}$. 
As before, we say that the velocity field $v$ \emph{induces} the probability path $p_t$ if the pair $(p_t,v)$ satisfies the continuity equation.

Flow matching methods do not parametrise the flow $\phi_t$ directly.
Instead, they learn a parametrised vector field $v_\theta(t,\z)$ whose
induced dynamics transport the base distribution to the target
distribution. Once the vector field is learned, the corresponding flow
$\phi_t$ is obtained by solving the associated ODE.

\paragraph{Training loss.} To learn the vector field, \cite{lipman2023flowmatching} proposes to train a parametric model $\v_\theta(x,t)$ using the flow-matching objective given by 
\begin{equation}
    \label{eq:FM_loss}    
    \cL_\text{FM}(\theta) \defeq \mathbb{E}_{t,p_t(x)}||\v_\theta(x,t) - \v_t(\x)||^2,
\end{equation}
where $\v_t(\x)$ is a vector field that induces the desired probability path. Observe that this is a supervised learning objective that assumes that we have access to the valid field $\v_t(\x)$, and that the parametric field $\v_\theta(x,t)$ is \emph{regressed} onto $\v_t(\x)$. This means that if $ \cL_\text{FM}(\theta) = 0$, then $\v_\theta(x,t)$ also generates the desired probability path. 

\begin{remark}[On the expectation in the flow-matching loss]
The expectation in eq.~\eqref{eq:FM_loss} reflects that the objective
is defined over the entire probability path. Rather than matching the
velocity field at a single point in time, the loss averages the
discrepancy between the parametric field $\v_\theta(\x,t)$ and the
target field $\v_t(\x)$ across both time and space, weighted by the
probability density $p_t(\x)$. Intuitively, this encourages the learned
vector field to reproduce the correct transport dynamics along the
whole path connecting the base and target distributions, while placing
more emphasis on regions where probability mass is concentrated.
\end{remark}

Though interpretable and simple, the FM training loss is rather naive: it is intractable in it own since neither $p_t$ nor $\v_t$ area available. In particular, there are a number of choices of probability paths such that $p_1(\x) = p(\x)$, and also several vector fields that can (perhaps redundantly) move mass according the source and target measures. The following approach takes a conditional perspective to circumvent the problems presented by the FM loss. 

\paragraph{Conditional probability paths.} Consider a particular data sample $\x_1$, and construct the conditional probability  path $p_t(\x \mid \x_1) $ with the boundary conditions\footnote{We start conditioning the path here to a terminal value $\x_1$, however, the conditioning statement can be much more flexible than that. We will revisit more expressive conditioning elements later in this section.}
\begin{align}
    t=0 : p_0(\x \mid \x_1) &= p(\x)\quad \leftarrow \quad \text{the source distribution}\\
    t=1 : p_1(\x \mid \x_1) &= \cN(\x \mid \x_1,\sigma^2\I) \quad \leftarrow \quad \text{with small } \sigma^2.
\end{align}
Conditional to the terminal sample, the evolution of the distribution above describes how samples from the source distribution flow towards a neighbourhood around $\x_1$ at $t=1$, rather than transporting mass directly toward the full data distribution. These trajectories induced by each datapoint can be then averaged over the conditioning endpoint $\x_1 \sim p(\x_1)$ to recover the marginal probability path. That is, 
\begin{equation}
    \label{eq:marginal_prob_path}
    p_t(\x) = \int p_t(\x \mid \x_1) q(\x_1)\d\x_1.
\end{equation}
The initial and terminal conditions for this marginal probability path are 
\begin{align}
    t=0 : p_0(\x) &=  \int p(\x) q(\x_1)\d\x_1 = p(\x)\\
    t=1 : p_1(\x) &=  \int p_t(\x \mid \x_1) q(\x_1)\d\x_1 =  \cN(\x \mid \x_1,\sigma^2\I) \ast q(\x),
\end{align}
meaning that if $\sigma \ll 1$, then $p_1(\x)\approx q(\x)$.

We now need to identify a corresponding vector field that produces the above probability path. \cite{lipman2023flowmatching} identifies the following result. 

\begin{theorem}[From conditional to marginal vector field]
    \label{thm:lipman_VF}
    The marginal vector field defined by 
    \begin{equation}
        \label{eq:marginal_vactor_field}
        \v_t(\x) = \int \v_t(\x \mid \x_1) 
        \frac{p_t(\x \mid \x_1) q(\x_1)}{p_t(\x)}\,\d\x_1, 
    \end{equation}
    generates the marginal probability path in 
    eq.~\eqref{eq:marginal_prob_path}.
\end{theorem}

\begin{proof}
The proof follows by showing that the marginal pair 
$ p_t(\x)$ and $\v_t(\x)$ also satisfy the continuity equation.

\begin{align*}
    \partial_t p_t(\x) 
    &= 
    \int \partial_t p_t(\x \mid \x_1)\, q(\x_1)\d\x_1 
    &&\text{(eq.~\eqref{eq:marginal_prob_path} and swap diff/integral)} \\[4pt]
    &= 
    - \int \nabla_\x \cdot 
    \big[ \v_t(\x \mid \x_1) p_t(\x \mid \x_1) \big] 
    q(\x_1)\d\x_1 
    &&\text{(conditional continuity equation)} \\[4pt]
    &= 
    - \nabla_\x \cdot 
    \left[\int 
    \v_t(\x \mid \x_1) p_t(\x \mid \x_1) 
    q(\x_1)\d\x_1 \right] 
    &&\text{(swap diff/integral)} \\[4pt]
    &= 
    - \nabla_\x \cdot 
    \left[
    p_t(\x)\!
    \int \v_t(\x \mid \x_1)
    \frac{p_t(\x \mid \x_1) q(\x_1)}{p_t(\x)}
    \d\x_1
    \right] 
    &&\text{(multiplying and dividing by $p_t(\x)$)} \\[4pt]
    &= 
    - \nabla_\x \cdot 
    \big[p_t(\x)\v_t(\x)\big]
    &&\text{(eq.~\eqref{eq:marginal_vactor_field})} .
\end{align*}

Hence $p_t(\x)$ and $\v_t(\x)$ satisfy the continuity equation,
which implies that $\v_t(\x)$ generates the marginal probability
path. The interchange of derivatives and integrals above is justified
under standard regularity assumptions (e.g., smooth vector fields and
integrable densities), ensuring the operations are well defined.
\end{proof}

\paragraph{Conditional flow matching.} Although Theorem \ref{thm:lipman_VF} provides a meaningful construction for the marginal vector field, this is again intractable as it depends on the data distribution $q(\cdot)$. However, we return to the flow-matching loss introduced above and extend a it to be implemented with conditional vector fields. This yields an alternative loss termed conditional flow matching defined by 

\begin{equation}
    \label{eq:CFM_loss}    
    \cL_\text{CFM}(\theta) \defeq \mathbb{E}_{t,q(\x_1), p_t(x \mid x_1)}||\v_\theta(x,t) - \v_t(\x \mid \x_1)||^2,
\end{equation}

Critically, this objective is tractable since it only requires sampling
from quantities that are directly available during training. In
particular, data points $\x_1 \sim q(\x_1)$ can be drawn from the
dataset, samples $\x \sim p_t(\x \mid \x_1)$ can be generated from the
chosen conditional probability path, and the corresponding conditional
vector field $\v_t(\x \mid \x_1)$ is analytically specified by the path
construction. As a result, the expectation in
eq.~\eqref{eq:CFM_loss} can be estimated using Monte Carlo samples
without requiring explicit knowledge of the marginal density $p_t(\x)$
or the marginal vector field $\v_t(\x)$.

The result linking this loss to the (marginal) FM loss is presented as follows 

\begin{theorem}
    If $p_t(\x)>0,\, \forall \x\in\R^D, t\in[0,1]$, then, $\cL_\text{FM}(\theta) = \cL_\text{CFM}(\theta), \forall \theta, $ up to constants that are independent of $\theta$. As a consequence, we have 
    \[
    \nabla_\theta \cL_\text{FM}(\theta) = \nabla_\theta\cL_\text{CFM}(\theta),
    \]
    therefore, the flow-matching model can be trained by stocastich gradient descent on the tractable $\cL_\text{CFM}(\theta)$.
\end{theorem}

\begin{proof}
    See \cite[Appendix]{lipman2023flowmatching}.
\end{proof}

\paragraph{Choices of the probability path and the associated vector field.} An alternative for the conditional probability path is as follows
\begin{equation}
    \label{eq:conditional_Gaussian_PP}
    p_t(\x \mid \x_1) = \cN(\x ; \mu_t(\x_1), \sigma_t^2(\x_1)\I), 
\end{equation}
with $\mu_t$ and $\sigma_t$ time-dependent functions. We can then set the boundary conditions: 
\begin{align}
    t=0 &:  \mu_0=0,\,\sigma_t= 1\, \Rightarrow p_0(\x) = \cN(\x \mid 0,\I)\\
    t=1 &: \mu_0=\x_1,\,\sigma_t= \sigma_\text{min}\, \Rightarrow p_(\x) = \cN(\x \mid \x_1,\sigma^2_\text{min}\I),
\end{align}
which means that the probability path (conditioned on an observation $\x_1$) goes from a standard normal on $\R^D$ to a Gaussian concentrated on $\x_1$. 

Among multiple possible flows that generate that probability path, we choose
\begin{equation}
    \label{eq:conditional_Gaussian_flow}
    \phi(\x) = \mu_t(\x_1) + \sigma_t(\x_1)\x; \quad \x\sim\cN(0,\I),
\end{equation}
where it is clear that eq.~\eqref{eq:conditional_Gaussian_flow} generates eq.~\eqref{eq:conditional_Gaussian_PP}. 


\begin{proposition}
    \label{prop:conditional_vector_field}
    The conditional vector field given by 
    \[
    \v_t (\x \mid \x_1) = \mu_t'(\x_1) + \frac{\sigma_t'(\x_1)}{\sigma_t(\x_1)}\left(\x  - \mu_t(\x_1)\right),
    \]
    generates the flow in eq.~\eqref{eq:conditional_Gaussian_flow}.
\end{proposition}
\begin{proof}
    Since the governing ODE is 
\(
\frac{\d}{\dt}\phi_t(\x) = \v_t(\phi_t(\x) \mid \x_1), 
\)
the proof proceeds by calculating the derivative of $\phi_t(\x)$ and expressing it as a function of the the same quantity. From eq.~\eqref{eq:conditional_Gaussian_flow} we know that 
\[
\x = \frac{\phi_t(\x) - \mu_t(\x_1)}{\sigma_t(\x_1)},
\]
combining this with the derivative of eq.~\eqref{eq:conditional_Gaussian_flow}, we have
\begin{align*}
    \frac{\d}{\dt}\phi_t(\x) 
    &= 
    \mu_t'(\x_1) + \sigma_t'(\x_1)\x,\\
    &= 
    \mu_t'(\x_1) + \sigma_t'(\x_1)\frac{\phi_t(\x) - \mu_t(\x_1)}{\sigma_t(\x_1)},
\end{align*}
thus concluding the proof.

    
\end{proof}

To construct a particular flow, let us consider the OT between the source $\cN(0,\I)$ and the (conditional) target $\cN(\x_1,\sigma^2_\text{min}\I)$. Recall that the geodesic interpolation is 
\begin{align*}
    \phi_t(\x) 
    &=  (1-t)\phi_0(\x) + t\phi_1(\x) \\
    &=  (1-t)\x + t(\x_1 + \x\sigma_\text{min}) \\ 
    &=  (1-t(1-\sigma_\text{min}))\x + t\x_1.
\end{align*}
This allows for identifying the time-indexed mean and variance as---see eq.~\eqref{eq:conditional_Gaussian_flow} 
\begin{align*}
    \mu_t(\x_1) &= t\x_1,\quad \text{going from $0$ to $\x_1$}\\
    \sigma_t(\x_1) &= 1-t(1-\sigma_\text{min}), \quad \text{going from $1$ to $\sigma_\text{min}$}.
\end{align*}
Lastly, evaluating these expressions into Prop.~\ref{prop:conditional_vector_field} gives the conditional vector field asºº
\begin{align}
\v_t(\x \mid \x_1)
&=
\frac{\x_1 - (1-\sigma_{\min})\,\x}
     {1 - t(1-\sigma_{\min})}.
\end{align}



